\documentclass[a4paper,12pt]{article}
\usepackage[T2A]{fontenc}
\usepackage[utf8]{inputenc}
\usepackage[russian]{babel}
\usepackage{amsmath,amssymb,mathtools}
\renewcommand{\familydefault}{\sfdefault}
\usepackage{icomma}
\usepackage[a4paper,margin=18mm,headheight=15pt]{geometry}
\usepackage{microtype}
\usepackage{tikz}
\usetikzlibrary{calc,arrows.meta,positioning,shapes.geometric}
\usepackage{tabularx,booktabs,longtable}
\usepackage{enumitem}
\usepackage{hyperref}
\usepackage{parskip}
\usepackage{fancyhdr}
\usepackage{setspace}
\usepackage{xcolor}

\definecolor{accent}{HTML}{2B6F73}
\definecolor{darkaccent}{HTML}{17494C}
\definecolor{textgray}{HTML}{2D3338}
\definecolor{softgray}{HTML}{6C747A}
\definecolor{pale}{HTML}{EEF5F4}

\hypersetup{colorlinks=true,linkcolor=darkaccent,urlcolor=darkaccent}
\setlength{\parindent}{0pt}
\onehalfspacing
\setlist[itemize]{leftmargin=1.5em,itemsep=0.3em,topsep=0.4em}
\setlist[enumerate]{leftmargin=1.8em,itemsep=0.45em,topsep=0.5em}
\pagestyle{fancy}
\fancyhf{}
\fancyhead[L]{\small\color{softgray}Неравенства с модулем}
\fancyhead[R]{\small\color{softgray}04.10.26}
\fancyfoot[C]{\small\color{softgray}\thepage}
\renewcommand{\headrulewidth}{0.2pt}
\renewcommand{\footrulewidth}{0pt}

\newcommand{\R}{\mathbb{R}}
\newcommand{\union}{\;\cup\;}
\newcommand{\lessonsection}[1]{%
  \vspace{0.5em}
  {\Large\bfseries\color{darkaccent}#1}\par\vspace{0.35em}
  {\color{accent}\rule{\linewidth}{0.6pt}}\vspace{0.6em}
}
\newcommand{\key}[1]{\textbf{\color{darkaccent}#1}}

\begin{document}

% ---------- Title ----------
\thispagestyle{empty}
\begin{center}
\vspace*{28mm}
{\Large\color{softgray}ЧЕК-ЛИСТ}\par
\vspace{8mm}
{\fontsize{27}{32}\selectfont\bfseries\color{darkaccent}Неравенства с модулем}\par
\vspace{4mm}
{\Large системы, совокупности, вложенные модули и замена}\par
\vspace{18mm}
{\large Анна Трапезникова}\par
\vspace{3mm}
{\normalsize 10 класс · углублённый уровень}\par
\vspace{16mm}
{\large 04.10.26}\par
\vfill
{\normalsize\color{textgray}Лёвин Артём Александрович эксклюзивно для Анны Трапезниковой}\par
\vspace{8mm}
\end{center}

\begin{tikzpicture}[remember picture,overlay]
  \draw[accent!70,line width=1.1pt]
    ($(current page.south west)+(12mm,16mm)$)
    .. controls ($(current page.south west)+(65mm,30mm)$)
    and ($(current page.south east)+(-70mm,3mm)$)
    .. ($(current page.south east)+(-12mm,19mm)$);
\end{tikzpicture}
\clearpage

% ---------- Main checklist ----------
\lessonsection{1. Главная идея}

Модуль переводит задачу в язык \key{расстояний} и почти всегда приводит к одному из двух логических объектов:
\begin{itemize}
  \item \key{система} --- должны выполняться все условия одновременно; на числовой прямой берём пересечение;
  \item \key{совокупность} --- достаточно выполнения хотя бы одного условия; на числовой прямой берём объединение.
\end{itemize}

Для выражений $F(x)$ и $G(x)$ удобно помнить четыре равносильных перехода:
\[
\begin{aligned}
|F(x)|<G(x)
&\iff
\begin{cases}
F(x)<G(x),\\
F(x)>-G(x),
\end{cases}
\\[1ex]
|F(x)|\le G(x)
&\iff
\begin{cases}
F(x)\le G(x),\\
F(x)\ge -G(x),
\end{cases}
\\[1ex]
|F(x)|>G(x)
&\iff
\left[
\begin{array}{l}
F(x)>G(x),\\
F(x)<-G(x),
\end{array}
\right.
\\[1ex]
|F(x)|\ge G(x)
&\iff
\left[
\begin{array}{l}
F(x)\ge G(x),\\
F(x)\le -G(x).
\end{array}
\right.
\end{aligned}
\]

\key{Проверка здравым смыслом.} Если справа стоит отрицательное число, некоторые ветви могут сразу стать тождественно истинными или невозможными. Это полезно замечать до вычислений.

\lessonsection{2. Двойное неравенство с модулем}

Двойное неравенство сначала разбивается на \key{систему} двух условий.

\textbf{Пример.} Решить
\[
1\le |x-2|<4.
\]
Записываем
\[
\begin{cases}
|x-2|\ge1,\\
|x-2|<4.
\end{cases}
\]
Первое условие даёт совокупность $x\le1$ или $x\ge3$, второе --- систему $-2<x<6$. Пересекаем результаты:
\[
\boxed{x\in(-2;1]\cup[3;6)}.
\]

\key{Практическое правило.} Для сложной конструкции рисуйте отдельную числовую прямую для каждой системы или совокупности. Так проще контролировать пересечения и объединения.

\lessonsection{3. Правая часть тоже содержит $x$}

Равносильные переходы работают и тогда, когда справа стоит выражение с переменной.

\textbf{Пример.}
\[
|-x+3|<x+5.
\]
Получаем систему
\[
\begin{cases}
-x+3<x+5,\\
-x+3>-(x+5).
\end{cases}
\]
Второе неравенство превращается в $3>-5$ и выполняется при любом $x$, поэтому остаётся
\[
-2x<2\quad\Longrightarrow\quad x>-1.
\]
Ответ: $(-1;+\infty)$.

\lessonsection{4. Вложенные модули}

Если модуль вложен в модуль, двигайтесь \key{снаружи внутрь}: сначала раскрывается внешний модуль, затем внутренний.

\textbf{Пример.}
\[
\bigl||x+5|-3\bigr|\ge4.
\]
Сначала
\[
\left[
\begin{array}{l}
|x+5|-3\ge4,\\
|x+5|-3\le-4.
\end{array}
\right.
\]
То есть
\[
\left[
\begin{array}{l}
|x+5|\ge7,\\
|x+5|\le-1.
\end{array}
\right.
\]
Вторая ветвь невозможна, потому что $|x+5|\ge0$. Первая даёт
\[
x\le-12\quad\text{или}\quad x\ge2.
\]

\lessonsection{5. Замена $t=|x|$}

Если один и тот же фрагмент $|x|$ повторяется несколько раз, удобно ввести
\[
t=|x|,\qquad t\ge0.
\]
Далее задача решается относительно $t$. Обратную замену делайте \key{после полного решения неравенства относительно $t$}.

\textbf{Пример.}
\[
\frac{|1-|x||}{1+|x|}\ge\frac12.
\]
Пусть $t=|x|\ge0$. Тогда $1+t\ge1>0$, поэтому можно умножать неравенство на $2(1+t)$ без изменения знака:
\[
2|1-t|\ge1+t.
\]
Отсюда
\[
t\in\left[0;\frac13\right]\cup[3;+\infty).
\]
Возвращаемся к $x$:
\[
\boxed{x\in(-\infty;-3]\cup\left[-\frac13;\frac13\right]\cup[3;+\infty)}.
\]

\key{Критическая проверка.} Умножать неравенство на выражение с переменной можно только тогда, когда его знак установлен. При умножении на отрицательное выражение знак неравенства меняется.

\clearpage

% ---------- Flowchart page ----------
\thispagestyle{plain}
\begin{center}
{\Large\bfseries\color{darkaccent}Алгоритм: как разбирать неравенство с модулем}\par
\vspace{8mm}
\end{center}

\begin{tikzpicture}[
  node distance=9mm,
  startstop/.style={ellipse,draw=accent,line width=0.9pt,align=center,minimum width=42mm,minimum height=10mm,fill=pale},
  process/.style={rectangle,rounded corners=2mm,draw=accent,line width=0.9pt,align=center,text width=118mm,minimum height=12mm},
  decision/.style={diamond,aspect=2.7,draw=accent,line width=0.9pt,align=center,text width=45mm,inner sep=1.5mm,fill=pale},
  arrow/.style={-{Stealth[length=2.4mm]},line width=0.9pt,color=darkaccent}
]
\node[startstop] (start) {Начало};
\node[process,below=of start] (prep) {Упростите выражение. Зафиксируйте ОДЗ и знаки множителей или знаменателей, если они важны.};
\node[decision,below=of prep] (repeat) {Повторяется фрагмент $|x|$?};
\node[process,below left=13mm and -7mm of repeat,text width=55mm] (subst) {Введите $t=|x|$, обязательно запишите $t\ge0$.};
\node[process,below right=13mm and -7mm of repeat,text width=55mm] (outer) {Работайте с внешним модулем напрямую.};
\node[process,below=25mm of repeat] (rule) {Для знаков $<$, $\le$ используйте систему; для $>$, $\ge$ --- совокупность. Решите полученные линейные или рациональные неравенства.};
\node[process,below=of rule] (combine) {На числовой прямой: система $\to$ пересечение, совокупность $\to$ объединение. Вложенные модули раскрывайте снаружи внутрь.};
\node[decision,below=of combine] (wasSub) {Была замена $t=|x|$?};
\node[process,below left=13mm and -7mm of wasSub,text width=55mm] (back) {Сделайте обратную замену и решите условия с $|x|$.};
\node[process,below right=13mm and -7mm of wasSub,text width=55mm] (check) {Проверьте границы, строгие знаки и логическую структуру ответа.};
\node[startstop,below=27mm of wasSub] (finish) {Ответ};

\draw[arrow] (start) -- (prep);
\draw[arrow] (prep) -- (repeat);
\draw[arrow] (repeat) -- node[above left,font=\small]{да} (subst);
\draw[arrow] (repeat) -- node[above right,font=\small]{нет} (outer);
\draw[arrow] (subst) |- (rule);
\draw[arrow] (outer) |- (rule);
\draw[arrow] (rule) -- (combine);
\draw[arrow] (combine) -- (wasSub);
\draw[arrow] (wasSub) -- node[above left,font=\small]{да} (back);
\draw[arrow] (wasSub) -- node[above right,font=\small]{нет} (check);
\draw[arrow] (back) |- (finish);
\draw[arrow] (check) |- (finish);
\end{tikzpicture}

\clearpage

% ---------- Mistakes ----------
\lessonsection{6. Частые ошибки}

\begin{enumerate}
  \item \key{Путаница системы и совокупности.} Система означает пересечение, совокупность --- объединение.
  \item \key{Автоматическое умножение на выражение с переменной.} Сначала установите его знак. Для $t=|x|\ge0$ знаменатель $1+t$ всегда положителен.
  \item \key{Игнорирование невозможной ветви.} Условие вида $|F(x)|<-2$ решений не имеет; условие $|F(x)|>-2$ выполняется при любом допустимом $x$.
  \item \key{Слишком ранняя обратная замена.} Сначала полностью решите задачу относительно $t$, учтите $t\ge0$, затем возвращайтесь к $x$.
  \item \key{Ошибка на последнем объединении.} После верных вычислений обязательно отдельно сформируйте итоговое множество решений.
\end{enumerate}

\lessonsection{7. Быстрая самопроверка}

Перед записью ответа проверьте:
\begin{itemize}
  \item соответствует ли знак модуля системе или совокупности;
  \item учтены ли строгие/нестрогие границы;
  \item нет ли ветви, которая всегда истинна или всегда ложна;
  \item установлен ли знак выражения, на которое вы умножали или делили;
  \item при замене $t=|x|$ использовано ли ограничение $t\ge0$;
  \item верно ли выполнены пересечение и объединение промежутков.
\end{itemize}

\clearpage

% ---------- Training ----------
\thispagestyle{plain}
\begin{center}
{\Large\bfseries\color{darkaccent}Тренировочный блок · Путь А}\par
\vspace{2mm}
{\color{softgray}10 заданий · от базовых переходов к комбинированным задачам}\par
\end{center}
\vspace{5mm}

\begin{enumerate}
  \item Решите неравенство
  \[
  |x-4|<3.
  \]

  \item Решите неравенство
  \[
  |2x+1|\ge5.
  \]

  \item Решите двойное неравенство
  \[
  1\le|x-2|<4.
  \]

  \item Решите неравенство
  \[
  |-x+3|<x+5.
  \]

  \item Решите неравенство
  \[
  |4x+5|\ge2x-7.
  \]

  \item Решите неравенство
  \[
  \bigl||x+5|-3\bigr|\ge4.
  \]

  \item Решите неравенство
  \[
  |3-|x||\ge2.
  \]

  \item Решите неравенство
  \[
  \frac{|1-|x||}{1+|x|}\ge\frac12.
  \]

  \item Решите неравенство
  \[
  \frac{(|x|-2)(|x|-5)}{|x|+1}<0.
  \]

  \item Решите неравенство
  \[
  \frac{|2-|x||}{1+|x|}\le\frac13.
  \]
\end{enumerate}

\clearpage

% ---------- Answers ----------
\thispagestyle{plain}
\begin{center}
{\Large\bfseries\color{darkaccent}Ответы}\par
\vspace{6mm}
\end{center}

\renewcommand{\arraystretch}{1.35}
\begin{tabularx}{\linewidth}{@{}cX@{}}
\toprule
\textbf{№} & \textbf{Ответ} \\
\midrule
1 & $x\in(1;7)$ \\
2 & $x\in(-\infty;-3]\cup[2;+\infty)$ \\
3 & $x\in(-2;1]\cup[3;6)$ \\
4 & $x\in(-1;+\infty)$ \\
5 & $x\in\R$ \\
6 & $x\in(-\infty;-12]\cup[2;+\infty)$ \\
7 & $x\in(-\infty;-5]\cup[-1;1]\cup[5;+\infty)$ \\
8 & $x\in(-\infty;-3]\cup\left[-\frac13;\frac13\right]\cup[3;+\infty)$ \\
9 & $x\in(-5;-2)\cup(2;5)$ \\
10 & $x\in\left[-\frac72;-\frac54\right]\cup\left[\frac54;\frac72\right]$ \\
\bottomrule
\end{tabularx}

\vfill
{\small\color{softgray}Лёвин Артём Александрович эксклюзивно для Анны Трапезниковой}\par

\end{document}
